% Revised \section{Conclusion and Discussion} for the IJOC manuscript.
% The existing first two paragraphs (framework summary + method summary)
% remain unchanged. The revision tightens the empirical-findings
% paragraph and the practical-implications paragraph to match the
% sharpened, data-backed story: DFL and a prediction-fairness penalty
% are the primary drivers of the regret-fairness trade-off, while MOO
% offers distinctive characteristics rather than uniform dominance.

% ----- REPLACE the paragraph that starts "We evaluate FDFL in a
% healthcare resource allocation problem using a synthetic dataset..."
% with the following: -----

We evaluate FDFL on two resource-allocation tasks: a healthcare
allocation problem derived from real patient records (48{,}784
patients, analytic decision gradients) and a synthetic multidimensional
knapsack (SPSA-based stochastic decision gradients). The two settings
are chosen to span the quality of the decision-gradient estimator,
which the empirical results identify as the main moderator of which
training strategy wins. On healthcare, where the analytical
vector-Jacobian product from Proposition~\ref{prop:analytical_gradient_group-1}
is available, every decision-focused end-to-end method beats every
two-stage baseline on regret at $\alpha = 2$ by $\approx 16$--$17\%$,
and the $\lambda$-sweep converts a clean fraction of that regret
budget into lower prediction-fairness violation---at MAD $\alpha = 0.5$,
for example, FDFL-Scal and FPLG reduce prediction-fairness violation
by about $75\%$ going from $\lambda = 0$ to $\lambda = 2$ at a regret
cost below $0.001$. Compared with DFL without prediction fairness,
FDFL thus attains notably better prediction fairness with minimal
loss in decision quality. On the knapsack, where the decision gradient
is approximated by SPSA, the end-to-end benefit persists for the
scalarized and $\mu$-anchored variants but becomes fragile for the
unanchored variant ($\mu = 0$) and for the multi-objective handlers,
which occasionally diverge on individual seeds. Adding any positive
prediction weight $\mu > 0$ restores stability, and the choice of
$\mu$ is second-order thereafter.

% ----- REPLACE the paragraph that starts "We highlight two practical
% implications of our study" with the following (keeps the same
% opening sentence but sharpens the two bullet-level claims): -----

We highlight two practical implications. First, E2EFO provides a
general and flexible framework for integrating fairness into
data-driven decision systems, and the distinction between prediction
fairness and decision fairness is essential: prediction fairness
corresponds to procedural equity---comparable predictive performance
across groups---while decision fairness captures outcome equity---equitable
distribution of allocation impacts. Our $\lambda$-sweep evidence shows
that these two forms of fairness are simultaneously achievable in the
FDFL framework without substantial regret loss, provided the decision
gradient is accurate enough to carry the regret signal on its own.
Second, the empirical results suggest a matching between the
decision-gradient regime and the training strategy. When the decision
gradient is analytically available, both the scalarized family and
the multi-objective handlers are viable; the multi-objective family
offers distinctive characteristics---most notably, no manual $\lambda$
tuning, and a strong $\alpha = 0.5$ advantage that comes from
resolving gradient conflict rather than from trading regret for
fairness---but does not uniformly dominate scalarization. When the
decision gradient is approximate, a positive fixed prediction weight
$\mu > 0$ is the first stabilizer to add; multi-objective projection
methods remain sensitive to gradient noise and require targeted
corrections (such as per-objective normalization, documented in
Appendix~\ref{app:mu-followup}) that resolve the directional imbalance
without eliminating the tail risk. This matching of training strategy
to gradient regime is itself a prescriptive contribution of the paper.
